\documentclass[11pt, a4paper]{article}

% --- UNIVERSAL PREAMBLE BLOCK ---
\usepackage[a4paper, top=2cm, bottom=2.5cm, left=2.5cm, right=2.5cm, marginparwidth=4cm]{geometry}
\usepackage{fontspec}
\usepackage[english, bidi=basic, provide=*]{babel}
\usepackage{multicol} 
\usepackage{xcolor}    
\usepackage{enumitem}
\usepackage{pdfpages}  

\babelprovide[import, onchar=ids fonts]{english}

% Use Noto Serif as the main font
\babelfont{rm}{Noto Serif}

% --- LAYOUT & UTILITIES ---
\usepackage{lineno}       
\usepackage{setspace}     
\usepackage{fancyhdr}     

% Line number settings
\modulolinenumbers[1]
\setlength\linenumbersep{2em}
\renewcommand\linenumberfont{\normalfont\tiny\sffamily\color{gray}}

% Header and Footer height adjustments
\setlength{\headheight}{14pt}
\setstretch{1.3} 

\pagestyle{fancy}
\fancyhf{}
\fancyhead[L]{\footnotesize \textsc{Reading Circle Handout}}
\fancyhead[R]{\footnotesize \textit{Du côté de chez Swann - Marcel Proust}}
\fancyfoot[R]{\footnotesize \textit{editor: minjun}}
\fancyfoot[C]{\footnotesize \thepage}
\fancyfoot[L]{\footnotesize \textit{2026 april 28}}
\renewcommand{\headrulewidth}{0pt} 

% --- END PREAMBLE ---

\begin{document}

% --- MINIMAL TITLE SECTION ---
\begin{flushleft}
    \vspace*{-1cm}
    {\Huge \textbf{Swann's Way }} \\ 
    \vspace{0.2em}
    {\large Marcel Proust (The Madeleine Episode) }\\ 
    
\end{flushleft}

% --- READING TASKS SECTION ---
\vspace{1em}
\section*{Reading Tasks: 20-Minute Focus}
\vspace{-0.5em}

\begin{itemize}[leftmargin=1.5em, itemsep=0.5em]
    \item \textbf{Sensory Triggers:} Circle specific words used to describe the physical sensation of the tea and cake.
    \item \textbf{The Process of Seeking:} Underline the narrator's transition from "mechanically" eating to the active, conscious "search" within his mind.
\end{itemize}
\vspace{0em}

\linenumbers 

\section*{I. THE SENSATION}

For many years already, everything about Combray that was not the theatre and drama of my bedtime had ceased to exist for me, when one day in winter, as I came home, my mother, seeing that I was cold, suggested that, contrary to my habit, I have a little tea. I refused at first and then, I do not know why, changed my mind. She sent for one of those squat, plump cakes called \textit{petites madeleines} that look as though they have been moulded in the grooved valve of a scallop-shell. And soon, mechanically, oppressed by the gloomy day and the prospect of a sad future, I carried to my lips a spoonful of the tea in which I had let soften a piece of madeleine.

But at the very instant when the mouthful of tea mixed with cake-crumbs touched my palate, I quivered, attentive to the extraordinary thing that was happening in me. A delicious pleasure had invaded me, isolated me, without my having any notion as to its cause. It had immediately made the \textbf{vicissitudes} of life unimportant to me, its disasters \textbf{innocuous}, its \textbf{brevity} illusory, acting in the same way that love acts, by filling me with a precious essence: or rather this essence was not in me, it was me. I had ceased to feel I was \textbf{mediocre}, \textbf{contingent}, mortal. Where could it have come to me from—this powerful joy?

\medskip
\noindent\textbf{Vocabulary:}\\
\textbf{vicissitudes} — the changes and setbacks of life; \textbf{innocuous} — harmless; \textbf{brevity} — shortness of life; \textbf{mediocre} — ordinary, insignificant; \textbf{contingent} — dependent on something else, not certain

\section*{II. THE SEARCH FOR TRUTH}

I sensed that it was connected to the taste of the tea and the cake, but that it went infinitely far beyond it, could not be of the same nature. Where did it come from? What did it mean? How could I grasp it? I drink a second mouthful, in which I find nothing more than in the first, a third that gives me a little less than the second. It is time for me to stop, the \textbf{virtue} of the drink seems to be diminishing. It is clear that the truth I am seeking is not in the drink, but in me. 

The drink has awoken it in me, but does not know that truth, and cannot do more than repeat indefinitely, with less and less force, this same testimony which I do not know how to interpret and which I want at least to be able to ask of it again and find again, intact, available to me, soon, for a decisive clarification. I put down the cup and turn to my mind. It is up to my mind to find the truth. But how? What grave uncertainty, whenever the mind feels overtaken by itself; when it, the seeker, is also the \textbf{obscure} country where it must seek and where all its baggage will be nothing to it. Seek? Not only that: create. It is face to face with something that does not yet exist and that only it can accomplish, then bring into its light.

\medskip
\noindent\textbf{Vocabulary:}\\
\textbf{virtue} — power, effect, or quality (archaic); \textbf{obscure} — unclear, hidden, mysterious

\section*{III. THE RESURFACING}

And I begin asking myself again what it could be, this unknown state which brought with it no logical proof, but only the evidence of its \textbf{felicity}, its reality, and in whose presence the other states of consciousness faded away. I want to try to make it reappear. I go back in my thoughts to the moment when I took the first spoonful of tea. I find the same state, without any new clarity. I ask my mind to make another effort, to bring back once more the sensation that is slipping away. 

Then for a second time I create an empty space before it, I confront it again with the still recent taste of that first mouthful and I feel something quiver in me, shift, try to rise, something that seems to have been \textbf{unanchored} at a great depth; I do not know what it is, but it comes up slowly; I feel the resistance and I hear the murmur of the distances \textbf{traversed}. 

Undoubtedly what is fluttering this way deep inside me must be the image, the visual memory which is attached to this taste and is trying to follow it to me. But it is struggling too far away, too confusedly... Will it reach the surface of my \textbf{limpid} consciousness—this memory, this old moment which the attraction of an identical moment has come so far to summon, to move, to raise up from my very depths? I don’t know.

\medskip
\noindent\textbf{Vocabulary:}\\
\textbf{felicity} — intense happiness, bliss; \textbf{unanchored} — detached, set adrift; \textbf{traversed} — crossed, traveled through; \textbf{limpid} — clear, transparent

\section*{IV. THE REVELATION}

And suddenly the memory appeared. That taste was the taste of the little piece of madeleine which on Sunday mornings at Combray, when I went to say good morning to her in her bedroom, my Aunt Léonie would give me after dipping it in her \textbf{infusion} of tea or \textbf{lime-blossom}. The sight of the little madeleine had not recalled anything to me before I tasted it; perhaps because, of these recollections abandoned so long outside my memory, nothing survived, everything had come apart...

But, when nothing \textbf{subsists} of an old past, after the death of people, after the destruction of things, alone, \textbf{frailer} but more enduring, more immaterial, more persistent, more faithful, smell and taste still remain for a long time, like souls, remembering, waiting, hoping, on the ruin of all the rest, bearing without giving way, on their almost \textbf{impalpable} droplet, the immense \textbf{edifice} of memory.

And as soon as I had recognized the taste of the piece of madeleine dipped in lime-blossom tea that my aunt used to give me... immediately the old grey house on the street, where her bedroom was, came like a stage-set to attach itself to the little wing opening on to the garden; and with the house the town, from morning to night and in all weathers, the Square, the streets where I went to do errands, the paths we took if the weather was fine. 

And as in that game in which the Japanese amuse themselves by filling a porcelain bowl with water and steeping in it little pieces of paper until then \textbf{indistinct}, which, the moment they are immersed in it, stretch and shape themselves, colour and differentiate, become flowers, houses, human figures, firm and recognizable, so now all the flowers in our garden and in M. Swann’s park, and the water-lilies on the Vivonne, and the good people of the village and their little dwellings and the church and all of Combray and its surroundings, all of this which is assuming form and substance, emerged, town and gardens alike, from my cup of tea.

\medskip
\noindent\textbf{Vocabulary:}\\
\textbf{infusion} — tea or herbal extract; \textbf{lime-blossom} — linden tree flower (tilleul); \textbf{subsists} — survives, remains; \textbf{frailer} — weaker, more fragile; \textbf{impalpable} — intangible, unable to be touched; \textbf{edifice} — a large building; a grand structure (here: the structure of memory); \textbf{indistinct} — unclear, blurry

\begin{center}
    \small \textsc{- fin -}
\end{center}

\newpage
\thispagestyle{empty}
\nolinenumbers
\newgeometry{top=1.2cm, bottom=1.2cm, left=1.5cm, right=1.5cm}

% --- APPENDIX SECTION ---
\vspace*{0.3em}
\section*{Appendix: The Proust Questionnaire}
\addcontentsline{toc}{section}{Appendix: The Proust Questionnaire}

\noindent\rule{\linewidth}{0.5pt}

\vspace{0.3em}
\noindent\textit{The Proust Questionnaire is a set of questions designed to reveal a person's true character. Made famous by French novelist Marcel Proust, it consists of 35 questions covering happiness, fears, loves, habits, and values. The origin of this questionnaire can be traced back to a 19th-century parlour game, but it was Proust's answers in 1886 (at age 14) that cemented its fame. Today, answering these questions remains a popular activity for understanding oneself and others.}

\vspace{0.3em}
\begin{multicols}{2}
\begin{enumerate}[leftmargin=0.5em, itemsep=0.3em, label=\textbf{\arabic*.}]
    \item What is your idea of perfect happiness?

    \hrulefill

    \item What is your greatest fear?

    \hrulefill

    \item What is the trait you most deplore in yourself?

    \hrulefill

    \item What is the trait you most deplore in others?

    \hrulefill

    \item Which living person do you most admire?

    \hrulefill

    \item What is your greatest extravagance?

    \hrulefill

    \item What is your current state of mind?

    \hrulefill

    \item What do you consider the most overrated virtue?

    \hrulefill

    \item On what occasion do you lie?

    \hrulefill
    \item What does it mean to you to suffer?

    \hrulefill
    \item Which living person do you most despise?

    \hrulefill
    \item What is the quality you most like in a man?

    \hrulefill
    \item What is the quality you most like in a woman?

    \hrulefill
    \item Which words or phrases do you most overuse?

    \hrulefill
    \item What or who is the greatest love of your life?

    \hrulefill
    \item When and where were you happiest?

    \hrulefill
    \item Which talent would you most like to have?

    \hrulefill
    \item If you could change one thing about yourself, what would it be?

    \hrulefill
    \item What do you consider your greatest achievement?

    \hrulefill
    \item If you were to die and come back as a person or thing, what would it be?

    \hrulefill
    \item Where would you most like to live?

    \hrulefill
    \item What is your most treasured possession?

    \hrulefill
    \item What do you regard as the lowest depth of misery?

    \hrulefill
    \item What is your favorite occupation?

    \hrulefill
    \item What is your most marked characteristic?

    \hrulefill
    \item What do you most value in your friends?

    \hrulefill
    \item Who are your favorite writers?

    \hrulefill
    \item Who is your favorite hero of fiction?

    \hrulefill
    \item Which historical figure do you most identify with?

    \hrulefill
    \item Who are your heroes in real life?

    \hrulefill
    \item What are your favorite names?

    \hrulefill
    \item What is it that you most dislike?

    \hrulefill
    \item What is your greatest regret?

    \hrulefill
    \item How would you like to die?

    \hrulefill
    \item What is your motto?

    \hrulefill
\end{enumerate}
\end{multicols}

\restoregeometry

\end{document}